\chapter{Collaudo, tipi di dato e input in Java}
\label{chap:collaudo-tipi-input-java}

Il \cref{chap:introduzione-java} ha introdotto la struttura di un programma
Java, le classi della libreria standard e la distinzione tra variabili
primitive e riferimenti a oggetti. Si approfondiscono ora l'uso delle
classi già disponibili, i tipi fondamentali e le operazioni sui loro
valori. Lo studio delle stringhe e dell'ingresso standard completa le
funzionalità necessarie per scrivere programmi il cui comportamento
dipende dai dati forniti dall'utente.\footnote{Materiale di riferimento:
S. Canazza, \emph{Collaudare una classe; tipi di dati e input}, slide
del corso di Fondamenti di Informatica,
\nolinkurl{2026_FI_08_canazza_Java_Collaudo_TipiDati_Input.pdf}.}

\section{Programmi di collaudo e pacchetti}
\label{sec:fi-java-collaudo}

L'uso di una classe non richiede di conoscerne l'implementazione:
è sufficiente sapere quali costruttori e metodi siano disponibili,
quali parametri accettino e quali risultati producano. Per verificare
il comportamento di una classe si può scrivere un \emph{programma di
collaudo}, costituito da una classe dotata di un metodo \texttt{main}.
Questo metodo costruisce gli oggetti necessari, ne invoca i metodi e
visualizza i risultati delle operazioni. Il confronto con i risultati
attesi permette di individuare eventuali comportamenti non corretti;
una singola esecuzione, tuttavia, non dimostra la correttezza generale
della classe.

Le classi della libreria standard sono organizzate in \emph{pacchetti}
(\emph{package}). La classe \texttt{Rectangle}, già incontrata
nella \cref{sec:fi-java-riferimenti}, appartiene a \texttt{java.awt}.
Per utilizzare il nome semplice della classe si introduce la dichiarazione
\begin{center}
  \lstinline|import java.awt.Rectangle;|
\end{center}
all'inizio del file sorgente, prima della classe che la utilizza.
In generale, la sintassi è \texttt{import nomePacchetto.NomeClasse;}.
L'importazione rende disponibile il nome della classe, non crea
oggetti né esegue il suo codice. Le classi del pacchetto
\texttt{java.lang}, tra cui \texttt{System}, \texttt{String},
\texttt{Math} e \texttt{Integer}, sono invece accessibili senza
importazioni esplicite.

Il programma seguente crea un rettangolo, ne modifica le coordinate
mediante \texttt{translate} e legge la nuova posizione:
\begin{unipdcode}[language=Java]
import java.awt.Rectangle;

public class MoveTester {
    public static void main(String[] args) {
        Rectangle box = new Rectangle(5, 10, 20, 30);
        box.translate(15, 25);
        System.out.println(box.getX());
        System.out.println(box.getY());
    }
}
\end{unipdcode}

\begin{unipdterminal}
20.0
35.0
\end{unipdterminal}
La traslazione somma \(15\) alla coordinata iniziale \(x=5\) e
\(25\) alla coordinata iniziale \(y=10\). Il metodo \texttt{getX()}
restituisce un valore \texttt{double}, perciò il risultato è
visualizzato con la parte decimale.

\section{Tipi fondamentali e rappresentazione numerica}
\label{sec:fi-java-tipi-fondamentali}

Le variabili di tipo primitivo conservano direttamente un valore,
mentre le variabili di tipo classe contengono un riferimento a un
oggetto, come spiegato nella \cref{sec:fi-java-riferimenti}.
Java definisce otto \emph{tipi fondamentali} o \emph{primitivi}:
quattro tipi interi, due in virgola mobile, \texttt{char} e
\texttt{boolean}. Le dimensioni e le rappresentazioni dei tipi
numerici e di \texttt{char} sono stabilite dal linguaggio; per
\texttt{boolean} il linguaggio specifica i valori ammissibili,
ma non una dimensione fissa in memoria.

\begin{table}[htbp]
  \centering
  \begin{tabular}{lll}
    \hline
    Tipo & Dimensione & Valori o caratteristiche \\
    \hline
    \texttt{byte} & 8 bit & interi da \(-128\) a \(127\) \\
    \texttt{short} & 16 bit & interi da \(-32768\) a \(32767\) \\
    \texttt{int} & 32 bit & interi da \(-2^{31}\) a \(2^{31}-1\) \\
    \texttt{long} & 64 bit & interi da \(-2^{63}\) a \(2^{63}-1\) \\
    \texttt{float} & 32 bit & virgola mobile, precisione singola \\
    \texttt{double} & 64 bit & virgola mobile, precisione doppia \\
    \texttt{char} & 16 bit & unità di codice UTF-16 \\
    \texttt{boolean} & non fissata & \texttt{true} oppure \texttt{false} \\
    \hline
  \end{tabular}
  \caption{Tipi primitivi di Java. La dimensione in memoria di
  \texttt{boolean} non è specificata come numero fisso di byte.}
  \label{tab:fi-java-tipi-primitivi}
\end{table}

I quattro tipi interi con segno usano il \emph{complemento a due},
introdotto nel \cref{chap:rappresentazione-informazioni}. Il tipo
\texttt{int} è adatto ai normali interi quando il suo intervallo
è sufficiente; per valori esterni a tale intervallo si usa
\texttt{long}, se la sua maggiore capacità è adeguata.
Un letterale intero senza suffisso è normalmente di tipo
\texttt{int}; se il valore non rientra nel suo intervallo,
è necessario indicare il tipo \texttt{long} mediante il suffisso
\texttt{L}, come in
\texttt{long popolazione = 3000000000L;}.

Le classi involucro della libreria standard espongono limiti
utilizzabili nei programmi: per esempio
\texttt{Integer.MIN\_VALUE}, \texttt{Integer.MAX\_VALUE},
\texttt{Long.MIN\_VALUE} e \texttt{Long.MAX\_VALUE}.
Per \texttt{float} e \texttt{double} le costanti
\texttt{MAX\_VALUE} individuano il massimo numero positivo finito;
\texttt{MIN\_VALUE} indica invece il minimo valore positivo
\emph{non nullo}, non il numero negativo più piccolo.

\subsection{Overflow e divisione per zero}
\label{sec:fi-java-overflow}

Se un'operazione fra \texttt{int} o \texttt{long} produce un valore
non rappresentabile, il risultato viene ricondotto all'intervallo
del tipo mediante \emph{wrap-around}, senza segnalare un'eccezione.
Per esempio, \texttt{Integer.MAX\_VALUE + 1} produce
\texttt{Integer.MIN\_VALUE}. Per ottenere un controllo esplicito,
la classe \texttt{Math} fornisce metodi come
\texttt{Math.addExact} e \texttt{Math.multiplyExact}, che lanciano
\texttt{ArithmeticException} in caso di overflow.

\begin{unipdcode}[language=Java]
int massimo = Integer.MAX_VALUE;
System.out.println(massimo + 1);
System.out.println(Math.addExact(100, 200));
\end{unipdcode}
\begin{unipdterminal}
-2147483648
300
\end{unipdterminal}

Tra le normali operazioni aritmetiche sugli interi, la divisione
o il resto con divisore zero lanciano \texttt{ArithmeticException}.
Per i tipi \texttt{float} e \texttt{double}, che seguono lo standard
IEEE 754, un overflow può invece produrre \( +\infty \) oppure
\( -\infty \), mentre la divisione \(0.0/0.0\) produce
\texttt{NaN} (\emph{Not a Number}). Java mette a disposizione
\texttt{Double.POSITIVE\_INFINITY},
\texttt{Double.NEGATIVE\_INFINITY} e \texttt{Double.NaN},
con le corrispondenti costanti della classe \texttt{Float}.

\subsection{Precisione dei numeri in virgola mobile}
\label{sec:fi-java-precisione}

I valori \texttt{float} e \texttt{double} occupano rispettivamente
32 e 64 bit. Un letterale decimale senza suffisso ha normalmente
tipo \texttt{double}; per esprimere un letterale \texttt{float}
si usa \texttt{f}, come in \texttt{float x = 30.2f;}.
La doppia precisione rappresenta approssimativamente 15--17 cifre
decimali significative, mentre la singola precisione ne offre
circa 6--9.

Come già mostrato nel \cref{chap:rappresentazione-informazioni},
molti numeri decimali non hanno una rappresentazione binaria finita.
L'arrotondamento può quindi modificare il risultato di un calcolo.
L'espressione \texttt{4.35 * 100}, per esempio, può produrre
\texttt{434.99999999999994} anziché \texttt{435.0}.
A grandezze dell'ordine di \(10^{17}\), la distanza tra due
valori \texttt{double} consecutivi è già sufficiente a perdere
piccole variazioni: sottrarre \(50\) da \(10^{17}\) non dà
necessariamente il risultato matematico esatto.

\section{Conversioni numeriche e arrotondamento}
\label{sec:fi-java-conversioni}

Nell'assegnamento il tipo del valore prodotto dall'espressione
deve essere compatibile con quello della variabile. Java ammette
molte conversioni numeriche \emph{implicite}, dette anche
\emph{widening}, tra cui il passaggio da \texttt{int} a
\texttt{long} o a \texttt{double}. Non tutte queste conversioni
preservano necessariamente il valore esatto: un \texttt{int} può
perdere precisione se convertito in \texttt{float}, e un
\texttt{long} sufficientemente grande può perderla anche se
convertito in \texttt{double}.

Quando una conversione può richiedere un restringimento del tipo,
come da \texttt{double} a \texttt{int}, è necessario un
\emph{cast} esplicito. La scrittura \texttt{(int) x} elimina
la parte frazionaria di un valore finito nell'intervallo
rappresentabile da \texttt{int}, troncando verso zero: il cast
non effettua un arrotondamento all'intero più vicino.

Il comportamento del cast e di \texttt{Math.round} è distinto
anche per un numero negativo:

\begin{unipdcode}[language=Java]
double valore = -2.95;
int troncato = (int) valore;
long arrotondato = Math.round(valore);
System.out.println(troncato);
System.out.println(arrotondato);
\end{unipdcode}
\begin{unipdterminal}
-2
-3
\end{unipdterminal}
\texttt{Math.round(double)} restituisce un \texttt{long},
mentre la versione che riceve un \texttt{float} restituisce
un \texttt{int}. In corrispondenza di valori esattamente a metà
tra due interi, \texttt{Math.round} sceglie quello verso
\(+\infty\); non applica quindi una regola simmetrica rispetto
allo zero. Il metodo è preferibile alla scorciatoia
\texttt{(int)(x + 0.5)}, che per valori negativi può dare
risultati errati.

Il cast applicato a \texttt{100 * 4.35} restituisce
\texttt{434}, perché il valore binario elaborato può risultare
lievemente inferiore a \texttt{435}; il troncamento elimina
quindi anche quella piccola differenza. È necessario distinguere
l'errore di rappresentazione dal comportamento del cast.

Per elaborare interi che superano i limiti di \texttt{long},
la libreria \texttt{java.math} fornisce la classe
\texttt{BigInteger}. La classe \texttt{BigDecimal} consente
invece calcoli decimali con precisione scelta dal programmatore,
utili quando la precisione binaria di \texttt{double} non è
adeguata. A differenza degli operatori fra tipi primitivi,
le operazioni su tali oggetti si eseguono mediante metodi.

La moltiplicazione di due interi arbitrariamente grandi
può essere eseguita mediante \texttt{BigInteger.multiply}:

\begin{unipdcode}[language=Java]
import java.math.BigInteger;

public class BigNumbers {
    public static void main(String[] args) {
        BigInteger a = new BigInteger("123456789");
        BigInteger b = new BigInteger("987654321");
        BigInteger c = a.multiply(b);
        System.out.println(c);
    }
}
\end{unipdcode}

\begin{unipdterminal}
121932631112635269
\end{unipdterminal}

\section{Costanti ed espressioni aritmetiche}
\label{sec:fi-java-costanti-aritmetica}

Una variabile dichiarata con la parola chiave \texttt{final}
può ricevere un valore soltanto una volta. Questo consente di
attribuire un nome significativo a una quantità che non deve
essere riassegnata durante l'esecuzione. La sintassi usuale è
\texttt{final nomeTipo NOME\_COSTANTE = espressione;}.
Per convenzione, i nomi di tali costanti sono scritti in
maiuscolo, con parole separate da un carattere di sottolineatura.

\begin{unipdcode}[language=Java]
final double EURO_PER_DOLLAR = 0.72;
double dollars = 2.35;
double euro = dollars * EURO_PER_DOLLAR;
\end{unipdcode}

Un nome esplicito elimina l'ambiguità di un \emph{numero magico}
e permette di modificare un valore utilizzato in più punti del
programma senza cercarne tutte le occorrenze.

In un'espressione Java, \texttt{*}, \texttt{/} e \texttt{\%}
hanno precedenza su \texttt{+} e \texttt{-}; le parentesi tonde
permettono di specificare l'ordine desiderato. Quando entrambi
gli operandi della divisione sono interi, il risultato è il
quoziente intero troncato verso zero. L'operatore \texttt{\%}
restituisce il resto della divisione intera.
Per esempio, \texttt{7 / 4} vale \texttt{1}, mentre
\texttt{7 \% 4} vale \texttt{3}. Se almeno uno degli operandi
è \texttt{double}, l'operazione di divisione viene effettuata
in virgola mobile, come in \texttt{7 / 4.0}, che vale
\texttt{1.75}.

Quoziente e resto possono essere usati per separare
la parte intera e quella frazionaria di un importo
espresso in centesimi:

\begin{unipdcode}[language=Java]
public class Coins5 {
    public static void main(String[] args) {
        double euro = 2.35;
        final int CENT_PER_EURO = 100;

        int centEuro = (int) Math.round(euro * CENT_PER_EURO);
        int intEuro = centEuro / CENT_PER_EURO;
        centEuro = centEuro % CENT_PER_EURO;
        
        System.out.print(intEuro + " euro e ");
        System.out.println(centEuro + " centesimi");
    }
}
\end{unipdcode}

\begin{unipdterminal}
2 euro e 35 centesimi
\end{unipdterminal}
Il valore in euro viene trasformato in un numero intero di
centesimi mediante \texttt{Math.round}; quoziente e resto della
divisione per \(100\) forniscono poi, rispettivamente, euro
interi e centesimi residui.

Quando una variabile viene aggiornata mediante un'operazione
aritmetica, Java offre gli operatori di assegnamento composti:
\texttt{x += y} corrisponde, nei casi numerici qui esaminati,
a \texttt{x = x + y}, e analogamente si usano \texttt{-=},
\texttt{*=} e \texttt{/=}. Gli operatori \texttt{++} e
\texttt{--} incrementano o decrementano una variabile numerica
di un'unità: come istruzioni autonome,
\texttt{counter++;} e \texttt{counter = counter + 1;}
hanno lo stesso effetto sul valore di \texttt{counter}.

\section{La classe \texttt{Math} e i metodi statici}
\label{sec:fi-java-math}

I metodi della classe \texttt{Math} permettono di eseguire
calcoli che non sono espressi da un singolo operatore aritmetico.
Tra quelli introdotti figurano \texttt{Math.pow(x, y)} per
l'elevamento a potenza, \texttt{Math.sqrt(x)} per la radice
quadrata, \texttt{Math.log(x)} per il logaritmo naturale e
\texttt{Math.sin(x)} per il seno di un angolo espresso in
radianti. La classe espone inoltre le costanti
\texttt{Math.PI} e \texttt{Math.E}.

Questi sono \emph{metodi statici}: si invocano attraverso il
nome della classe, nella forma
\begin{center}
  \texttt{NomeClasse.nomeMetodo(argomenti)}.
\end{center}
Non è quindi necessario creare un oggetto della classe
\texttt{Math}. Un metodo di istanza, come
\texttt{System.out.println(...)} o
\texttt{box.translate(...)}, viene invece invocato mediante
un riferimento all'oggetto ricevente, secondo la
\cref{sec:fi-java-variabili}. La dichiarazione
\texttt{public static long round(double a)} indica, per
esempio, un metodo pubblico statico che riceve un
\texttt{double} e restituisce un \texttt{long}.

\section{Stringhe e conversioni testuali}
\label{sec:fi-java-stringhe-operazioni}

Una \texttt{String} rappresenta una sequenza di testo e, a
differenza dei tipi primitivi, è un oggetto. La classe è già
stata introdotta nella \cref{sec:fi-java-metodi-string},
dove si sono esaminati \texttt{length}, \texttt{replace} e
\texttt{toUpperCase}. Si considerano ora l'estrazione di
sottostringhe, la concatenazione e la conversione tra testo
e numeri.

\subsection{Sottostringhe e indici}
\label{sec:fi-java-substring}

Il metodo \texttt{substring(inizio, fine)} restituisce la
parte di una stringa che va dall'indice \texttt{inizio},
\emph{incluso}, all'indice \texttt{fine}, \emph{escluso}.
Gli indici iniziano da zero. Per una stringa di lunghezza
\(n\) sono ammissibili \(0\leq\texttt{inizio}\leq
\texttt{fine}\leq n\); la sottostringa ha lunghezza
\(\texttt{fine}-\texttt{inizio}\). Con un solo argomento,
\texttt{substring(inizio)} estrae i caratteri fino al termine
della stringa.

Le due forme del metodo \texttt{substring} possono essere
utilizzate sulla stessa stringa:

\begin{unipdcode}[language=Java]
String s = "Hello, World!";
String prima = s.substring(0, 5);
String seconda = s.substring(7);
System.out.println(prima);
System.out.println(seconda);
\end{unipdcode}
\begin{unipdterminal}
Hello
World!
\end{unipdterminal}
La prima chiamata estrae gli indici da \(0\) a \(4\), la
seconda quelli a partire dall'indice \(7\). Un indice
negativo, un indice finale superiore alla lunghezza o un
intervallo con inizio maggiore della fine provoca
\texttt{StringIndexOutOfBoundsException}. Se l'eccezione non
viene gestita, l'esecuzione termina e le istruzioni successive
non vengono eseguite.

L'espressione \texttt{s.substring(s.length() - 4)} estrae
gli ultimi quattro elementi della stringa, purché essa abbia
almeno quattro unità di lunghezza. Quando occorre limitare
un indice superiore alla lunghezza si può usare
\texttt{Math.min(s.length(), limite)}; ciò non elimina la
necessità di controllare anche gli altri vincoli sugli indici.

\subsection{Concatenazione e immutabilità}
\label{sec:fi-java-concatenazione}

Se uno degli operandi di \texttt{+} è una \texttt{String},
Java esegue una \emph{concatenazione}, convertendo l'altro
operando in testo se necessario. Non vengono aggiunti spazi
automaticamente: \texttt{15 + "euro"} produce
\texttt{"15euro"}, mentre \texttt{15 + " " + "euro"}
produce \texttt{"15 euro"}. Il metodo \texttt{concat}
permette anch'esso di unire stringhe.

Le stringhe Java sono \emph{immutabili}: i metodi come
\texttt{toUpperCase}, \texttt{toLowerCase}, \texttt{replace}
e \texttt{substring} restituiscono un risultato senza alterare
l'oggetto originale. L'assegnamento a una variabile
\texttt{String} può cambiarne il riferimento, ma non modifica
l'oggetto stringa a cui si riferiva in precedenza.

\begin{unipdcode}[language=Java]
String s = "Hello";
String trasformata = s.toUpperCase() + s.toLowerCase();
System.out.println(trasformata);
System.out.println(s);
\end{unipdcode}
\begin{unipdterminal}
HELLOhello
Hello
\end{unipdterminal}

\subsection{Conversione tra stringhe e numeri}
\label{sec:fi-java-parsing}

Una stringa che contiene cifre non diventa automaticamente un
numero: l'assegnamento \texttt{int eta = "28";} è
incompatibile con il tipo richiesto. Per ottenere il valore
numerico si utilizzano i metodi statici
\texttt{Integer.parseInt} e \texttt{Double.parseDouble}.
Se la stringa non rappresenta un valore valido nel formato
atteso, questi metodi lanciano
\texttt{NumberFormatException}.

\begin{unipdcode}[language=Java]
String testo = "28";
int eta = Integer.parseInt(testo);
double valore = Double.parseDouble("18.3");
String risultato = Integer.toString(eta);
\end{unipdcode}

Per la conversione inversa si possono usare
\texttt{Integer.toString}, \texttt{Double.toString} oppure
la concatenazione con una stringa vuota. Le classi
\texttt{Integer} e \texttt{Double} sono \emph{classi involucro}
(\emph{wrapper}) dei rispettivi tipi primitivi.

\section{Caratteri e sequenze di escape}
\label{sec:fi-java-char-escape}

Il tipo primitivo \texttt{char} conserva una singola unità
di codice UTF-16; i letterali di carattere sono racchiusi
tra apici singoli, come \texttt{'x'}, anziché tra virgolette
doppie. Rispetto a una stringa di lunghezza unitaria, una
variabile \texttt{char} rappresenta direttamente il valore
primitivo, senza richiedere un oggetto \texttt{String}.
Il metodo \texttt{charAt(i)} della classe
\texttt{String} restituisce l'unità di codice alla posizione
\(i\), usando la stessa indicizzazione a partire da zero
di \texttt{substring}. Analogamente, \texttt{length()} conta
le unità di codice UTF-16: un punto di codice Unicode rappresentato
da una coppia di surrogati occupa due posizioni, secondo
la \cref{sec:fi-caratteri}.
Se l'indice è fuori intervallo viene lanciata
\texttt{StringIndexOutOfBounds\allowbreak Exception}.

\begin{unipdcode}[language=Java]
String nome = "John";
char lettera = nome.charAt(2);
System.out.println(lettera);
for (int i = 0; i < nome.length(); i++) {
    System.out.print(nome.charAt(i));
}
System.out.println();
\end{unipdcode}
\begin{unipdterminal}
h
John
\end{unipdterminal}

Il ciclo \texttt{for} inizializza \texttt{i} a zero,
prosegue finché \texttt{i} è minore della lunghezza della
stringa e incrementa l\'indice dopo ogni iterazione.
In questo modo si visitano tutti gli elementi in ordine.

Gli operatori aritmetici possono promuovere i valori \texttt{char}
a \texttt{int}: \texttt{'a' + 'b'} produce \(97+98=195\),
non la concatenazione \texttt{"ab"}. Quest'ultima si ottiene
usando una stringa, per esempio
\texttt{"" + 'a' + 'b'}.

All'interno di stringhe e caratteri, alcune sequenze
introdotte da una barra rovesciata rappresentano caratteri
speciali. Tra le più comuni vi sono \lstinline|\n| per una
nuova riga, \lstinline|\t| per una tabulazione,
\lstinline|\r| per il ritorno carrello,
\lstinline|\b| per \emph{backspace},
\lstinline|\f| per \emph{form feed},
\lstinline|\"| per una virgoletta doppia e
\lstinline|\\| per la barra rovesciata stessa.
Una sequenza Unicode della forma \lstinline|\u00E9| codifica,
per esempio, la lettera \emph{é}.

\begin{unipdcode}[language=Java]
System.out.println("Hello, \"World\"!");
System.out.println("*\n**\n***");
\end{unipdcode}
\begin{unipdterminal}
Hello, "World"!
*
**
***
\end{unipdterminal}

La corretta visualizzazione di caratteri non ASCII dipende
anche dalla codifica e dalla configurazione del terminale.
Questo non significa che Java sia limitato ad ASCII: le
stringhe del linguaggio utilizzano Unicode, mentre lo
scambio di testo con l'ambiente esterno richiede una codifica
compatibile. I caratteri di controllo ASCII da 0 a 31,
richiamati nella \cref{sec:fi-caratteri}, comprendono anche
ETX ed EOT. Nei terminali, combinazioni come \texttt{Ctrl+C}
e \texttt{Ctrl+D} sono comunemente interpretate come richieste
di interruzione o di fine dell'ingresso, secondo le convenzioni
dell'ambiente; non sono normali caratteri da visualizzare.

\section{Formattazione dell'output}
\label{sec:fi-java-printf}

Il metodo \texttt{println} utilizza una rappresentazione
predefinita dei valori. Per scegliere il numero di cifre
decimali e allineare i risultati si può usare
\texttt{System.out.printf}. Il primo argomento è una
\emph{stringa di formato}, che contiene testo ordinario e
\emph{specificatori} introdotti da \texttt{\%}; gli argomenti
successivi forniscono i valori da visualizzare.

\begin{table}[htbp]
  \centering
  \begin{tabular}{ll}
    \hline
    Specificatore & Significato \\
    \hline
    \texttt{\%d} & intero decimale \\
    \texttt{\%x}, \texttt{\%o} & intero esadecimale, ottale \\
    \texttt{\%f} & numero in virgola mobile in notazione decimale \\
    \texttt{\%e}, \texttt{\%g} & notazione esponenziale, formato generale \\
    \texttt{\%s} & stringa o rappresentazione testuale \\
    \texttt{\%n} & terminatore di riga della piattaforma \\
    \hline
  \end{tabular}
  \caption{Alcuni specificatori utilizzabili con \texttt{printf}.}
  \label{tab:fi-java-formati}
\end{table}

In \lstinline|%5.2f| il numero \(5\) indica una larghezza
minima di cinque caratteri e \(2\) il numero di cifre dopo
il separatore decimale. Il valore è arrotondato ai fini
della visualizzazione; il contenuto della variabile non viene
modificato. Il separatore decimale può dipendere dalla
\emph{localizzazione} impostata nell'ambiente di esecuzione.
Usando \texttt{Locale.US} si ottiene il punto come separatore.

Tra i modificatori più comuni figurano \lstinline|-| per
l'allineamento a sinistra, \lstinline|0| per il riempimento
con zeri iniziali, \lstinline|+| per visualizzare anche il
segno positivo, \lstinline|(| per racchiudere tra parentesi
i valori negativi e \lstinline|,| per il raggruppamento
delle migliaia secondo la localizzazione. Per richiedere
lettere maiuscole si usa, quando disponibile, la conversione
maiuscola, ad esempio \lstinline|%E| invece di
\lstinline|%e|.

La stringa di formato permette di visualizzare importi
con due cifre decimali e colonne allineate:

\begin{unipdcode}[language=Java]
import java.util.Locale;

public class FormatTester {
    public static void main(String[] args) {
        double total = 3.50;
        final double TAX_RATE = 8.5;
        double tax = total * TAX_RATE / 100;
        System.out.printf(Locale.US, "Total:%5.2f%n", total);
        System.out.printf(Locale.US, "Tax:  %5.2f%n", tax);
    }
}
\end{unipdcode}

\begin{unipdterminal}
Total: 3.50
Tax:   0.30
\end{unipdterminal}
Il formato \lstinline|%5.2f| occupa almeno cinque posizioni,
mentre \lstinline|%n| inserisce un terminatore di riga
adatto alla piattaforma; in particolare, non è necessario
presupporre che tutti i sistemi usino il solo carattere
\lstinline|\n|.

\section{L'ingresso standard con \texttt{Scanner}}
\label{sec:fi-java-scanner}

Un programma che inizializza sempre gli stessi valori produce
il medesimo calcolo a ogni esecuzione. Per ricevere dati
variabili dall'utente, Java mette a disposizione il flusso
di ingresso standard \texttt{System.in}, un oggetto di tipo
\texttt{InputStream}. La classe \texttt{Scanner}, definita nel
pacchetto \texttt{java.util}, offre metodi per interpretare
come numeri o stringhe i dati letti da tale flusso.
\texttt{Scanner} può essere collegato anche ad altre sorgenti
di input, come un file; qui si considera l'ingresso standard.

Per creare uno \texttt{Scanner} collegato all'ingresso standard
si utilizza \texttt{new Scanner(System.in)}, dopo aver
importato \texttt{java.util.Scanner}. \texttt{System.in} e
lo \texttt{Scanner} sono oggetti distinti: il secondo legge
e interpreta i dati forniti dal primo. È opportuno chiudere
lo \texttt{Scanner} quando non serve più; chiudere uno
\texttt{Scanner} associato a \texttt{System.in} chiude però
anche il flusso sottostante.

\subsection{Lettura di numeri}
\label{sec:fi-java-scanner-numeri}

Il metodo \texttt{nextInt()} legge un valore di tipo
\texttt{int}, mentre \texttt{nextDouble()} legge un valore
di tipo \texttt{double}. Quando l'ingresso non contiene ancora
un valore completo, la chiamata può restare in attesa dei dati
forniti dall'utente. In caso di input non compatibile con
il tipo richiesto, \texttt{Scanner} può lanciare
\texttt{InputMismatchException}. Per uniformare l'ingresso
decimale alla convenzione con il punto si usa
\texttt{in.useLocale(Locale.US)}, importando anche
\texttt{java.util.Locale}.

Un programma può quindi acquisire le quantità durante
l'esecuzione e convertire il valore complessivo in euro:

\begin{unipdcode}[language=Java]
import java.util.Locale;
import java.util.Scanner;

public class Coins6 {
    public static void main(String[] args) {
        Scanner in = new Scanner(System.in);
        in.useLocale(Locale.US);
        System.out.println("Quante lire?");
        int lit = in.nextInt();
        System.out.println("Quanti euro?");
        double euro = in.nextDouble();
        System.out.printf(Locale.US, "Valore totale in euro %.2f%n",
                          euro + lit / 1936.27);
        in.close();
    }
}
\end{unipdcode}

Con gli input \texttt{15000} e \texttt{2.35}, forniti alle
due richieste, il programma stampa un totale arrotondato a
due cifre decimali. I valori sono letti durante l'esecuzione:
per cambiare l'importo non occorre modificare o ricompilare
il sorgente.

\subsection{Lettura di parole e righe}
\label{sec:fi-java-scanner-testo}

Il metodo \texttt{next()} restituisce il successivo
\emph{token} separato da spaziature; \texttt{nextLine()}
restituisce invece il resto della riga corrente, fino al
terminatore di riga, che non viene incluso nel risultato.
Di conseguenza, \texttt{next()} permette di leggere una
parola anche quando più valori sono sulla stessa riga,
mentre \texttt{nextLine()} consente di acquisire anche
stringhe contenenti spazi.

Quando una chiamata a \texttt{nextInt()} o \texttt{next()}
è seguita da \texttt{nextLine()}, il terminatore della riga
può essere ancora presente nell'ingresso. In tal caso
\texttt{nextLine()} può restituire una stringa vuota;
per leggere la riga successiva occorre prima consumare
quella corrente. Per acquisire una serie di dati interamente
per righe, è anche possibile usare sempre \texttt{nextLine()}
e convertire le righe numeriche con \texttt{Integer.parseInt}
o \texttt{Double.parseDouble}, come nella
\cref{sec:fi-java-parsing}.

Lettura per righe e conversione numerica si combinano nel
programma seguente, che costruisce una stringa a partire
dai dati forniti dall'utente:

\begin{unipdcode}[language=Java]
import java.util.Scanner;

public class MakePassword2 {
    public static void main(String[] args) {
        Scanner in = new Scanner(System.in);
        String firstName = in.nextLine();
        String lastName = in.nextLine();
        int age = Integer.parseInt(in.nextLine());
        in.close();
        String initials = firstName.substring(0, 1)
                        + lastName.substring(0, 1);
        System.out.println(initials.toLowerCase() + age);
    }
}
\end{unipdcode}

Se l'utente inserisce \texttt{Sergio}, \texttt{Canazza} e
\texttt{28} su tre righe, il programma combina le iniziali
minuscole con il valore numerico e stampa \texttt{sc28}.
La regola è ripresa esclusivamente come esempio di
manipolazione di stringhe: il risultato è prevedibile e
non deve essere utilizzato come password reale.

I metodi \texttt{next()} e \texttt{nextInt()} possono anche
leggere dati separati da spazi o distribuiti su righe
successive, poiché entrambi ignorano le spaziature iniziali.
La scelta tra lettura per token e per riga dipende quindi
dal formato previsto per i dati in ingresso, non soltanto
dal loro tipo.
